\documentclass[12pt]{article}
\usepackage[utf8]{inputenc}
\usepackage[T1]{fontenc}
\usepackage{hyperref}
\title{Cameroon Pidgin Language Resources for K-12 Curriculum Development}
\author{Compiled by Claude Sonnet 5.5 (claude-sonnet-5-5) for Andr\'as Kornai}
\date{}
\begin{document}
\maketitle

\section{Introduction}
Cameroon Pidgin English (ISO 639-3 \texttt{wes}) is an English-based creole and a widely used language of wider communication in Cameroon. Search results state that it has no standardised orthography, though it has a strong oral tradition and is used in broadcast media.

\section{Primary Educational Resources}
I found no school textbooks or primers in search results.

\section{Linguistic and Reference Materials}
\begin{itemize}
\item Overview: \url{https://en.wikipedia.org/wiki/Cameroonian_Pidgin_English}
\item Jean-Paul Kouega, \emph{A Dictionary of Cameroon Pidgin English Usage} (2008 and 2015 editions are mentioned). The Glottolog reference record I saw is: \url{https://glottolog.org/resource/reference/id/52263}
\item Ayafor and Green, \emph{Cameroon Pidgin English: A comprehensive grammar}, chapter listing: \url{https://lwc1.benjamins.com/catalog/loall.20.c12}
\item Survey of the linguistic preferences of Cameroon Pidgin English speakers (orthography and lexical choice): \url{https://benjamins.com/catalog/jpcl.9.1.04wil}
\item Spoken corpus of Cameroon Pidgin English, project description: \url{https://blog.philsoc.org.uk/2019/01/25/a-spoken-corpus-of-cameroon-pidgin-english-compilation-applications-and-next-steps/}
\item PHOIBLE source record (Mbah and Anderson 2005): \url{https://phoible.org/sources/1632_mbah_pny}
\end{itemize}

\section{Religious and Community Resources}
Search results indicate that the Catholic Church made the earliest attempt to codify the language (catechisms) and that Baptist missionaries adopted it for evangelism. A Bible Society of Cameroon New Testament exists. Its text appears to be the one excerpted in the Ayafor and Green grammar as ``Good Nyus fo Ol Pipol''. A scholarly chapter treats the translation as language elaboration (Muehleisen and Anchimbe); I saw no direct URL for that chapter, so it is not linked. Dramatized audio New Testament (Bible Society of Cameroon, 2000): \url{https://megavoice.com/media-cloud/m0b6334-new-testament-dramatized-wes-pidgin-cameroon-audio-bible/}

\section{Digital and Audio Resources}
\begin{itemize}
\item Global Recordings Network, Cameroon Pidgin: \url{https://globalrecordings.net/en/language/wes}
\item Dramatized audio New Testament (see above).
\end{itemize}

\section{Curriculum Development Context}
The language is spoken by millions, but orthography is not standardised and I found no school programme. A curriculum effort would need to settle a spelling standard, building on the Bible translation and the Kouega dictionary, with English and French as the official-language context.
\end{document}